% AI-produced by GPT-6 Astra, OpenAI. Reconstructed 20 November 2026.
% Unnumbered method lemma; pre-outage proof audited by root and
% singularities_duality. This is a self-contained finite affine proof.
\subsection*{Unnumbered finite-format extension across an affine folded disk}

\textbf{Lemma.}\quad
Let $K$ be a finite triangulation of a closed two-dimensional disk,
with every vertex on its boundary. Let
$\iota:|K|\to\mathbb R^N$ be affine on each triangle, and let
$T\subset\mathbb R^N$ be a nondegenerate affine triangle. Suppose
$g:\partial|K|\to\partial T$ is a PL homeomorphism, affine on each
boundary edge of $K$, and that the three corners of $T$ are images
of vertices of $K$. Suppose also that
\[
 \|\iota(v)-g(v)\|_2\le\delta
 \quad\hbox{for every vertex }v\hbox{ of }K.
 \tag{FD.1}
\]
For every $\varepsilon>0$ there is a PL homeomorphism
$h:|K|\to T$, equal to $g$ on the entire boundary, with
\[
 \sup_{x\in|K|}\|\iota(x)-h(x)\|_2<\delta+\varepsilon.
 \tag{FD.2}
\]
For fixed combinatorial data, these homeomorphisms admit one finite
semialgebraic formula with a positive scalar perturbation parameter.
The number of PL pieces needed is bounded in terms of the given
triangulation, independently of $\varepsilon$.

\begin{proof}
\textbf{The initial affine map and its cost.}\quad
Write the boundary vertices in the cyclic order induced from the
positive orientation of $T$ as $v_1,\ldots,v_n$, and put $p_i=g(v_i)$.
They are distinct, occur in that cyclic order around $\partial T$,
and include all three corners. Let $q_0:|K|\to T$ be the affine
vertex map $v_i\mapsto p_i$ on every triangle. It need not be
injective: some triangles can collapse or overlap. Nevertheless,
for a point with barycentric coordinates $\lambda_i$ in any one
triangle, convexity of the norm and (FD.1) give
\[
 \|\iota(x)-q_0(x)\|_2
 \le\sum_i\lambda_i\|\iota(v_i)-p_i\|_2\le\delta.
 \tag{FD.3}
\]
The proof will replace $q_0$ by a nearby homeomorphism while retaining
its exact boundary map.

\textbf{A small strictly convex outward perturbation.}\quad
Fix an orientation of the affine plane of $T$. On each oriented side
$A\to B$ of $T$, choose a nonzero outward normal vector $n_{AB}$
in that plane. For a boundary vertex on this side, write
$p_i=A+\lambda_i(B-A)$, $0\le\lambda_i\le1$, and set
\[
 w_i=\lambda_i(1-\lambda_i)n_{AB},\qquad
 p_i(\tau)=p_i+\tau w_i.
 \tag{FD.4}
\]
At a corner the value is zero from either incident side, so this
definition is consistent. All corners remain fixed.

For three distinct points on the same side, with
$\lambda_i<\lambda_j<\lambda_k$, direct expansion gives
\[
 \begin{split}
 &\det\bigl(p_j(\tau)-p_i(\tau),p_k(\tau)-p_i(\tau)\bigr)\\
 &\quad=-\tau\det(B-A,n_{AB})
 (\lambda_j-\lambda_i)(\lambda_k-\lambda_i)(\lambda_k-\lambda_j)>0.
 \end{split}
 \tag{FD.5}
\]
Here $\det(B-A,n_{AB})<0$ because the triangle is positively oriented
and the normal is outward. Any three distinct boundary points not
lying on one side have strictly positive determinant in their cyclic
order already at $\tau=0$. There are finitely many such triples,
so all their determinants remain positive for sufficiently small
$\tau>0$. Equation (FD.5) handles every formerly collinear triple.
Thus, for some $\tau_0>0$, the points $p_i(\tau)$, $0<\tau<\tau_0$,
are the vertices of a strictly convex polygon $Q_\tau$ in the same
cyclic order. Indeed, every cyclic triple has positive determinant,
which gives the strict supporting-half-plane condition for every
polygon edge. Since its vertex set includes all three corners of
$T$, convexity also gives $T\subset Q_\tau$.

Put $M=\max_i\|w_i\|_2$. The boundary displacement is at most
$\tau M$. All choices above are finite semialgebraic ones; in
particular the common positivity condition determining an admissible
$\tau_0$ is a finite list of polynomial inequalities.

\textbf{Why the same disk triangulation straightens injectively.}\quad
A triangulated disk with all vertices on its boundary is an
outerplanar triangulation. We give the needed elementary argument.
If it has $n$ boundary vertices, Euler's formula and incidence
counting give $n-2$ triangles and $n-3$ interior edges. Its triangle
dual graph is connected and hence is a tree. A leaf supplies an ear
triangle with two boundary edges. Removing its ear vertex and that
triangle leaves a triangulated disk on the remaining cyclic vertices.
The assertion that the affine vertex map to a strictly convex
$n$-gon is a homeomorphism now follows by induction on $n$.
For $n=3$ it is an affine triangle homeomorphism. At an ear, the
corresponding convex ear triangle and the convex polygon on the
remaining vertices meet exactly in their common diagonal and cover
the original convex polygon. Paste the induction homeomorphism and
the affine ear map along that diagonal.

Apply this argument to the points $p_i(\tau)$. The affine vertex map
\[
 q_\tau:|K|\longrightarrow Q_\tau,
 \qquad q_\tau(v_i)=p_i(\tau),
 \tag{FD.6}
\]
is a PL homeomorphism. Barycentric interpolation gives
\[
 \sup_{x\in|K|}\|q_\tau(x)-q_0(x)\|_2\le\tau M.
 \tag{FD.7}
\]
This planar ear argument is the reason for the two-dimensional and
all-boundary-vertices hypotheses.

\textbf{Restoring the exact boundary.}\quad
Let $c$ be the centroid of the nondegenerate triangle $T$.
It lies in the interior of $T$ and therefore in the interior of
$Q_\tau$. The triangles
$[c,p_i(\tau),p_{i+1}(\tau)]$ form a fan triangulation of $Q_\tau$.
The triangles $[c,p_i,p_{i+1}]$ form a fan triangulation of $T$.
All target fan triangles are nondegenerate: consecutive boundary
points are distinct and their supporting side line does not contain
$c$. Map corresponding fan triangles affinely, sending $c$ to $c$
and $p_i(\tau)$ to $p_i$. This yields a PL homeomorphism
\[
 \beta_\tau:Q_\tau\longrightarrow T.
 \tag{FD.8}
\]
Its displacement on each fan triangle is a convex combination of
zero and the vectors $-\tau w_i$, so
\[
 \sup_{y\in Q_\tau}\|\beta_\tau(y)-y\|_2\le\tau M.
 \tag{FD.9}
\]
The composite $h_\tau=\beta_\tau q_\tau$ equals $g$ on every entire
boundary edge, not just at the vertices, since all three maps are
affine there with the required endpoint values. Equations
(FD.3), (FD.7), and (FD.9) give
\[
 \sup_{x\in|K|}\|\iota(x)-h_\tau(x)\|_2
 \le\delta+2\tau M.
 \tag{FD.10}
\]
Choose
\[
 0<\tau<\min\{\tau_0,\varepsilon/[2(1+M)]\}.
 \tag{FD.11}
\]
Then (FD.2) follows. This also covers $M=0$, when there are no
noncorner boundary vertices and no perturbation is needed.

\textbf{Finite formula and complexity.}\quad
The maps $q_\tau$ are affine on the fixed triangles of $K$.
The fan maps $\beta_\tau$ are given by finitely many barycentric
formulas; their inverse matrices have nonzero determinants for all
admissible $\tau$. Thus their evaluation and inverse evaluation,
and those of $h_\tau$, are described by one finite semialgebraic
formula in $\tau$ and the fixed vertex data. Selection of an
admissible $\tau$ for a given positive error is also semialgebraic.
No additional increasingly fine mesh is used as the error tends
to zero.

More concretely, intersect the image of each of the $n-2$ triangles
of $K$ under $q_\tau$ with each of the $n$ fan triangles in $Q_\tau$.
An intersection of two plane triangles is a convex polygon with at
most six vertices. It can be triangulated with at most four triangles;
lower-dimensional intersections need no two-dimensional pieces.
Pulling these subdivisions back makes the composite affine on at
most $4n(n-2)$ two-dimensional pieces. The actual combinatorial
subdivision may change with the scalar parameter, but this finite
bound does not. This proves the asserted finite-format conclusion.
\end{proof}

\paragraph{Original-source context and scope record.}
This proof is self-contained. Its affine scaffold motivation was
checked in the following original full texts: A. Nayyeri and H. Xu,
\emph{On the Decidability of the Fr\'echet Distance between Surfaces},
author version dated 13 July 2017, 18 pages,
\url{https://web.engr.oregonstate.edu/~nayyeria/pubs/frechet.pdf},
Lemma~4, p.~7, and Lemma~5, pp.~7--8; and their
\emph{On Computing the Fr\'echet Distance Between Surfaces},
SoCG 2016, DOI \url{https://doi.org/10.4230/LIPIcs.SoCG.2016.55},
Lemma~6 and Lemma~7, p.~55:8, with the following shortcutting lemmas
on pp.~55:8--55:9. Both complete originals were reopened after
recovery and the exact stated portions reread.

The lemma supplies an exact-boundary finite-format extension for
this affine two-dimensional cell problem, with arbitrarily small
positive error slack. It does not produce a compatible bounded
scaffold for a curved definable surface family. It also does not
extend the convex straightening claim to higher-dimensional ball
triangulations. No general local-bundle-locus or homeomorphic graph
matching theorem is inferred from it.
